\section{The Chern character}
\label{sec:chernchar}

In this section we define the Chern character map $\ch\colon \K_0(X)
\ra \HH_0(X)$.  We discuss the relationship between our construction
and the one of Markarian~\cite[Definition~2]{Mar}.  Then we show that
the Chern character maps the Euler pairing to the Mukai pairing: we
call this the Semi-Hirzebruch-Riemann-Roch Theorem.

\subsection{Definition of the Chern character.}
Suppose $X$ is a space, and $\cE$ is an object in $\D(X)$. Consider
$\cE$ as an object of $\D(\pt\times X)$, i.e., as a kernel $\pt \ra
X$, so there is an induced linear map
\[ \cE_*\colon\HH_\blob(\pt) \ra \HH_\blob(X). \]
Now, because the Serre functor on a point is trivial, $\HH_0(\pt)$ is
canonically identifiable with $\Hom_{\pt\times \pt}(\cO_\pt, \cO_\pt)$
so there is a distinguished class $\bone \in \HH_0(\pt)$ corresponding
to the identity map.  Define the Chern character of $\cE$ as
\[ \ch(\cE) := \cE_*(\bone) \in \HH_0(X). \]
Graphically this has the following description:
\[\ch(\cE):=\pic{chE}.\]
Naturality of push-forward leads to the next theorem.

\begin{Theorem}
\label{thm:commut}
If $X$ and $Y$ are spaces and $\Phi\colon X\to Y$ is a kernel then the
diagram below commutes.
\[
\begin{diagram}[height=2em,width=2em,labelstyle=\scriptstyle]
\D(X) & \rTo^{\Phi\circ {-}} & \D(Y) \\
\dTo_{\ch} & & \dTo_{\ch} \\
\HH_0(X) & \rTo^{\Phi_*} & \HH_0(Y).
\end{diagram}
\]
\end{Theorem}

\begin{Proof}
  Let $\cE$ be an object of $\D(X)$.  We will regard it either as an
  object in $\D(X)$, or as a kernel $\pt \ra X$, and similarly we will
  regard $\Phi\circ\cE$ either as an object in $\D(Y)$ or as a kernel
  $\pt \ra Y$.  By Theorem~\ref{thm:functorial} we have
  \begin{align*}
    \Phi_* \ch(\cE) & = \Phi_* \left(\cE_* (\bone)\right) = (\Phi\circ
    \cE)_* \bone = \ch(\Phi\circ\cE).\qedhere
  \end{align*}
\end{Proof}

 \subsection{The Chern character as a map on K-theory.}
\label{subsec:iotae}
To show that the Chern character descends to a map on K-theory we give
a characterization of the Chern character similar to that of
Markarian~\cite{Mar}.

For any object $\cE\in \D(X)$, which is to be considered an object of
$\D( \pt\times X)$, there are the following two maps:
\[\iota_\cE\colon \HH_\blob(X)\to \Hom_{\D(X\times X)}^\blob(\cE,\SerK_X\circ \cE);
\qquad\pic{vinHH}\,\mapsto \,\pic{iEofvinHH}\]
and
\[\iota^\cE\colon \Hom_{\D(X\times X)}^\blob(\cE,\cE)\to\HH_\blob(X);
\qquad\pic{einHomEE}\mapsto \,\pic{homologyiEofeinHomEE}. \]

Recall that the Mukai pairing is a non-degenerate pairing on $\HH_\blob(X)$ and that the Serre pairing is a perfect pairing between $\Hom_{\D(X\times X)}^\blob(\cE,\SerK_X\circ \cE)$
and $\Hom_{\D(X\times X)}^\blob(\cE,\cE)$.
%The vector space $\HH_\blob(X)$ is self dual by virtue of the Mukai
%pairing. Similarly, $\Hom_{\D(X\times X)}^\blob(\cE,\SerK_X\circ \cE)$
%and $\Hom_{\D(X\times X)}^\blob(\cE,\cE)$ are dual by the Serre
%pairing. 
 With respect to these pairings the two maps $\iota^\cE$ and
$\iota_\cE$ are adjoint in the following sense.

\begin{Proposition}\label{prop:iotasadjoint}
  For $\phi\in \Hom^\blob(\cE,\cE)$ and $v\in \HH_\blob(X)$ the
  following equality holds:
  \[\langle v,\, \iota^\cE\phi\rangle_{\rm M} =\SP{\iota_\cE v}{\phi}\]
\end{Proposition}

\begin{Proof}
  Here in the third equality we use the invariance of the Serre trace
  under the partial trace map.
  \begin{align*}
    \MP{v}{\iota^\cE\phi}&=
    \MPbig{\pic{vinHH}}{\pic{homologyiEofeinHomEE}}
    :=\Tr\left(\pic{iEadjoint3}\right)\\&=
    \Tr\left(\pic{iEadjoint4}\right)
    =\Tr\left(\pic{iEadjoint5}\right)\\
    &=: \SPbig{\pic{iEofvinHH}\quad}{\quad\pic{einHomEE}} =:\SP{\iota_\cE
      v}{\phi}.\qedhere
  \end{align*}
\end{Proof}
\noindent
Note that, using this, the Chern character could have been defined as
\[\ch(\cE):=\iota^\cE(\id_\cE).\]
Then from the above proposition the following is immediate.

\begin{Lemma}
  For $v\in \HH_0(X)$ and $\cE\in \D(X)$ there is the equality
  \[\MP{v}{\ch(\cE)}=\Tr(\iota_\cE(v)),\]
  and this defines $\ch(\cE)$ uniquely.
\end{Lemma}

The fact that the Chern character descends to a function on the
K-group can now be demonstrated.

\begin{Proposition}
  For $\cE\in \D(X)$ the Chern character $\ch(\cE)$ depends only on
  the class of $\cE$ in $\K_0(X)$.  Thus the Chern character can be
  considered as a map
  \[ \ch\colon \K_0(X) \ra \HH_0(X). \]
\end{Proposition}

\begin{Proof}
  It suffices to show that if $\cF\ra \cG\ra \cH\ra \cF[1]$ is an
  exact triangle in $\D(X)$, then
  \[ \ch(\cF) - \ch(\cG) + \ch(\cH) = 0 \]
  in $\HH_0(X)$.

  For $\cG,\cH\in \D(X)$, $\alpha \colon \cG\to \cH$ and $v\in
  \HH_\blob(X)$ the diagram on the left commutes as it expresses the
  equality on the right:
  \[
  \begin{diagram}[height=2em,width=2em]
    \cG & \rTo^{\alpha} & \cH\\
    \dTo^{\iota_\cG( v)} & &
    \dTo_{\iota_\cH( v)}  \\
    \SerK_X \circ\cG & \rTo^{\Id_{\SerK_X}\circ \alpha} & \SerK_X
    \circ\cH.
  \end{diagram}\qquad;\qquad\pic{ChernOnK1}\ =\ \pic{ChernOnK2}\ .
  \]
  In other words, from an element $v\in \HH_\blob(X)$ we get
  $\tau_R(v)\in \Hom(\Id_X,\SerK_X)$, which in turn gives rise to
  a natural transformation between the functors $\Id_{\D(X)}\colon \D(X)\to
  \D(X)$ and $\SerK_X\circ{-}\colon \D(X)\to \D(X)$.  This leads to a
  map of triangles
  \[
  \begin{diagram}[height=2em,width=2em]
    \cF & \rTo & \cG & \rTo & \cH & \rTo & \cF[1] \\
    \dTo_{\iota_\cF( v)} & & \dTo_{\iota_\cG( v)} & &
    \dTo_{\iota_\cH( v)} & & \dTo_{\iota_\cF( v)[1]} \\
    S_X \cF & \rTo & S_X \cG & \rTo & S_X \cH & \rTo & S_X \cF[1].
  \end{diagram}
  \]
  Observe that if we represent the morphism $v$ by an actual map of
  complexes of injectives, and the objects $\cF$, $\cG$ and $\cH$ by
  complexes of locally free sheaves, then the resulting maps in the
  above diagram commute on the nose (no further injective or locally
  free resolutions are needed), so we can apply
  \cite[Theorem~1.9]{May} to get
  \[ \Tr_X(\iota_\cF( v)) - \Tr_X(\iota_\cG( v)) + \Tr_X(\iota_\cH(
  v)) = 0. \] Therefore, by the lemma above, for any $ v\in
  \HH_\blob(X)$,
  \[ \MP{v}{\ch(\cF)-\ch(\cG)+\ch(\cH)}= 0 .\]
  Since the Mukai pairing on $\HH_0(X)$ is
  non-degenerate, we conclude that
  \[ \ch(\cF) - \ch(\cG) + \ch(\cH) = 0. \qedhere\]
\end{Proof}

 \subsection{The Chern character and inner products.}
One reading of the Hirzebruch-Riemann-Roch Theorem is that it says
that the usual Chern character map $\ch\colon \K_0\to \H^\blob(X)$ is
a map of inner product spaces when $\K_0(X)$ is equipped with the
Euler pairing (see below) and $\H^\blob(X)$ is equipped with the
pairing $\langle x_1,x_2\rangle:=(x_1\cup x_2 \cup \Td_X)\cap [X]$.
It is shown in~\cite{CalHH2} that the Hochschild homology Chern
character composed with the Hochschild-Kostant-Rosenberg map
$I_{\text{HKR}}$ gives the usual Chern character:
\[\K_0\xrightarrow{\ch}\HH_0(X)\xrightarrow{I_\text{HKR}}
\bigoplus_p \H^{p,p}(X).\]
Here we show that the Hochschild homology Chern character is an
inner-product map when $\HH_\blob(X)$ is equipped with the Mukai
pairing.

First recall that the Euler pairing on $\K_0(X)$ is defined by
\[\chi(\cE,\cF):=\sum_i (-1)^i \dim \Ext^i_X(\cE, \cF). \]

\begin{Theorem}[Semi-Hirzebruch-Riemann-Roch]
\label{thm:hrr}
The Chern character
\[ \ch\colon \K_0\to \HH_0(X) \]
is a map of inner-product spaces: in other words, for $\cE, \cF\in
\D(X)$ we have
\[ \MP{\ch(\cE)}{\ch(\cF)} = \chi(\cE, \cF). \]
\end{Theorem}

\begin{Proof}
  The first thing to do is to get an interpretation of the Euler
  pairing.  Considering the totality of $\Ext$-groups
  $\Ext^\blob(\cE,\cF)$ as a graded vector space, the Euler
  characteristic is just the graded dimension of
  $\Ext^\blob(\cE,\cF)$, which is to say it is the trace of the
  identity map on $\Ext^\blob(\cE,\cF)$.  Moreover, if ${\mathbf R}\pi_*\colon
  \D(X)\to \D(\pt)$ is the derived functor coming from the map $X\to
  \pt$, then
  $$\Ext^\blob(\cE,\cF)\cong
  \H^\blob(X,\cE^\chk\otimes\cF)\cong {\mathbf R}\pi_*(\cE^\chk\otimes\cF)$$ and
  the latter is just the composition $\cE^\chk\circ \cF$, where
  $\cE^\chk$ and $\cF$ are considered as kernels respectively $X\to
  \pt$ and $\pt\to X$.  Thus, using the invariance of the Serre trace
  under the partial trace map,
  \begin{align*}
    \chi(\cE,\cF)&=\Tr(\Id_{\cE^\chk\circ \cF})
    =\Tr\left(\pic{HRR1}\right)
    =\Tr\left(\pic{HRR2}\right)\\&=\MPbig{\pic{chE}}{\pic{chF}}
    =\MP{\ch(\cE)}{\ch(\cF)}.\qedhere
  \end{align*}
\end{Proof}

 \subsection{Example.}  To have a non-commutative example at
hand, consider the case when $G$ is a finite group acting trivially
on a point.  The orbifold $BG$ is defined to be the global quotient $[\scdot/G]$ and then the category of coherent sheaves on the
orbifold $BG$ is precisely the category of finite
dimensional representations of $G$.  One can naturally identify
$\HH_0(BG)$ with the space of conjugation invariant functions on $G$,
and the Chern character of a representation $\rho$ is precisely the
representation-theoretic character of $\rho$.
See~\cite{WilGr} for details.
